\section{Script Listings}\label{app:scripts}

\iffullversion\else
Table~\ref{tab:assume} lists the deployment assumptions of
\S\ref{sec:scripts}.
% The deployment-assumptions table, mirroring spec 4.1's [A-1] to [A-6].
% Input in section 5 of the full version and in Appendix C of the short
% one, so that both cite it as Table~\ref{tab:assume}.
\begin{table}[tbp]
\centering
\caption{Deployment assumptions (spec 4.1, \texttt{spec/layertree.md}). A proposal \emph{defines} its opcode for one tapleaf version; a tier may \emph{assume} it in another.}
\label{tab:assume}
\setlength{\tabcolsep}{2pt}\scriptsize
\begin{tabular}{@{}lp{8.9cm}l@{}}
\toprule
 & assumption & tiers \\
\midrule
{[A-1]} & \op{TX}~\cite{optx} (a draft on a personal fork, defined for \texttt{0xC2} with BIP-440 costs) pushes amounts as 8-byte little-endian strings, which FULL compares with BIP-441 numerics & FULL, MID, CCV \\
{[A-2]} & CSFS (BIP-348, defined for \texttt{0xC0}) available in \texttt{0xC2} at the cost of one signature check & FULL, CCV \\
{[A-3]} & the EC-ops draft~\cite{ecops} available in \texttt{0xC2} at its stated costs & FULL \\
{[A-4]} & opcode numbers: \op{TX} and \op{EC\_POINT\_NEGATE} both claim \texttt{0xbd}, BIP-443 and \op{EC\_POINT\_ADD} both \texttt{0xbb}; the simulator parks \op{TX}, \op{TWEAKADD} and CCV at \texttt{0xbf}, \texttt{0xc0}, \texttt{0xc1}, so the golden scripts are byte-exact under that numbering only (sizes are unaffected) & all \\
{[A-5]} & \op{TX} in \texttt{0xC0} with \texttt{CScriptNum} outputs, which no proposal defines & MID, CCV \\
{[A-6]} & BIP-443 (defined for \texttt{0xC0}) also in \texttt{0xC2}; its sigop cost, open upstream, charged as one signature check; amount rules as of \texttt{bitcoin/bips} commit \texttt{24e96e87}. Conservation rests on the tier's \op{TX} equalities, so a change to BIP-443 can only make honest spends fail & CCV \\
{[MIN]} & \op{CHECKSIG} under key $G$ with $s = e + 1$ accepts only if the script-built message is the real one, barring a collision in the TapSighash or challenge hash & MIN \\
{[C]} & contention: $\kappa$ and $\beta$ of \S\ref{sec:liveness} (policy and fees, not consensus) & all \\
\bottomrule
\end{tabular}
\end{table}

\fi

Figure~\ref{fig:step} states the step every transition is built from; the listings below are taken from the reference generators and are what
the size figures of \S\ref{sec:eval} measure. Stack comments show the
main stack left to right with the top on the right, and the alt stack
after a bar.

\begin{figure}[h]
\centering
\begin{pcvstack}[center]
\scalebox{0.9}{\procedure[linenumbering]{$\mathsf{Step}(\pi, n_{\mathrm{old}}, n_{\mathrm{new}}, (r, S))$}{
  (r_0, S_0) \gets \fold(n_{\mathrm{old}}, \pi) \pccomment{$D-h$ unrolled levels} \\
  \pcassert r_0 = r \wedge S_0 = S \pccomment{\tr{CHECK\_OLD}: $S$ is input 0's amount} \\
  (r', S') \gets \fold(n_{\mathrm{new}}, \pi) \pccomment{same siblings, same sums} \\
  \pcreturn (r', S')
}}
\pcvspace
\scalebox{0.9}{\procedure[linenumbering]{$\mathsf{Final}((r', S'), m)$}{
  \pcif (r', S') = \Emp_D \pcthen \pcreturn \pccomment{terminal: output 0 free} \\
  (a_0, \spk_0) \gets \op{TX}(\text{output } 0) \\
  \pcassert \spk_0 = \TreeSpk(r', m) \wedge a_0 = S'
}}
\end{pcvstack}
\caption{The Merkle-sum step shared by all transitions, and the output rule. A transition is a sequence of steps, each proving its old node against the root the previous step produced (\tr{CHECK\_PREV}), followed by \textsf{Final}. Authorisation and payout checks are specific to each transition.}
\label{fig:step}
\end{figure}

\paragraph{One level of the fold (FULL tier, 88 bytes).}
\tr{LEVEL} advances both accumulators, the old node $(oh, os)$ and the
new node $(nh, ns)$, by one level using the same direction bit and
sibling $sc = h \,\|\, \encs(s)$ from the witness. $T$ is the
64-byte tag prefix of $\Hnode$; \textsf{PAD8} is
$\langle 2^{64}\rangle\ \op{ADD}$, which yields $\encs$ of the sum by
Lemma~\ref{lem:enc}.

\begin{lstlisting}
OP_6 OP_ROLL OP_6 OP_ROLL             # T oh os nh ns dir sc
OP_SIZE 41 OP_EQUALVERIFY             #   sc is exactly hash || enc64(sum)
OP_DUP 32 OP_8 OP_SUBSTR              # ... dir sc ss8   (sibling sum, raw le64)
OP_DUP OP_TOALTSTACK                  #                          | ss8
OP_3 OP_ROLL OP_TUCK OP_ADD           # T oh os nh dir sc ns ns'
OP_TOALTSTACK PAD8                    # T oh os nh dir sc ns9    | ss8 ns'
OP_3 OP_ROLL OP_SWAP OP_CAT           # T oh os dir sc nc        (nc = nh || enc64(ns))
OP_OVER OP_SWAP OP_3 OP_PICK          # T oh os dir sc sc nc dir
OP_NOTIF OP_SWAP OP_ENDIF OP_CAT      # T oh os dir sc pre_n     (dir ? sc||nc : nc||sc)
OP_5 OP_PICK OP_SWAP OP_CAT OP_SHA256 # T oh os dir sc nh'
OP_TOALTSTACK                         #                          | ss8 ns' nh'
OP_FROMALTSTACK OP_FROMALTSTACK OP_FROMALTSTACK
                                      # T oh os dir sc nh' ns' ss8
OP_5 OP_ROLL OP_TUCK OP_ADD           # T oh dir sc nh' ns' os os'
OP_TOALTSTACK PAD8                    # T oh dir sc nh' ns' os9  | os'
OP_5 OP_ROLL OP_SWAP OP_CAT           # T dir sc nh' ns' oc      (oc = oh || enc64(os))
OP_3 OP_ROLL OP_SWAP OP_4 OP_ROLL     # T nh' ns' sc oc dir
OP_NOTIF OP_SWAP OP_ENDIF OP_CAT      # T nh' ns' pre_o
OP_3 OP_PICK OP_SWAP OP_CAT OP_SHA256 # T nh' ns' oh'
OP_FROMALTSTACK OP_2SWAP              # T oh' os' nh' ns'
\end{lstlisting}

\paragraph{Checking the old root and writing the new one.}
After $D$ levels, \tr{CHECK\_OLD} compares the old accumulator with the
committed state: the sum against input~0's amount and the hash against
the witness root $r$ parked on the alt stack (itself bound to the input
scriptPubKey by the prologue). \tr{FINAL} constrains output~0 unless the
new root is the empty tree.

\begin{lstlisting}
# CHECK_OLD:  [T oh os nh ns] | ... r  ->  [T nh ns]
OP_2SWAP <SEL_IN_AMOUNT> OP_TX OP_NUMEQUALVERIFY OP_FROMALTSTACK OP_EQUALVERIFY

# FINAL:  [T nh ns] | m  ->  [1]
OP_ROT OP_DROP                                   # nh ns
OP_OVER <E_D.hash> OP_EQUAL
OP_IF   OP_2DROP OP_FROMALTSTACK OP_DROP OP_1    # terminal: output 0 free
OP_ELSE OP_SWAP OP_FROMALTSTACK TREE_SPK         # ns spk'
        OP_0 <SEL_OUT_SINGLE> OP_TX              # ns spk' amt0 spk0
        OP_ROT OP_EQUALVERIFY OP_NUMEQUAL        # 1
OP_ENDIF
\end{lstlisting}

\paragraph{The MID gadget.}
Under \texttt{0xC0} a computed 4-byte number $x$ is widened to 8 bytes by
two witness elements $z, w$: the script checks
$\mathsf{size}(x) + \mathsf{size}(z) = 4$ and $z \,\|\, w = \mathtt{00000000}$
and uses $x \,\|\, z \,\|\, w$ as $\lesf(x)$. Two such gadgets per level
and the per-field size checks account for almost all of MID's premium
over FULL.

\paragraph{The MIN message.}
The MIN exit assembles the BIP-341 signature message for input~0 with
hash type \texttt{SIGHASH\_DEFAULT} and checks it as described in
\S\ref{sec:min}. Fields marked \emph{derived} are computed by the script
from values it has verified; fields marked \emph{raw} are witness
elements the script does not constrain but must be presented
consistently, since the whole message is hashed and compared with the
real one through \op{CHECKSIG}.

\begin{center}\scriptsize
\begin{tabular}{@{}p{2.4cm}p{5.4cm}p{3.8cm}@{}}
\toprule
field & source & note \\
\midrule
epoch, hash type & constant $\mathtt{00\,00}$ & \texttt{SIGHASH\_DEFAULT} \\
version, lock time & raw & free \\
sha\_prevouts & raw & free \\
sha\_amounts & derived: $\mathsf{SHA256}(\lesf(S) \,\|\, \mathit{amt}_1)$ & $S$ verified by \tr{CHECK\_OLD}; $\mathit{amt}_1$ raw; two inputs exactly \\
sha\_scriptpubkeys & derived: $\mathsf{SHA256}(\mathtt{22} \,\|\, \TreeSpk(r,m) \,\|\, \mathit{spk}_1)$ & $\TreeSpk$ recomputed; $\mathit{spk}_1$ raw \\
sha\_sequences & raw & the field the spender grinds \\
sha\_outputs & derived: $\mathsf{SHA256}(\mathit{out}_0 \,\|\, \mathit{out}_1 \,\|\, \mathit{extra})$ & $\mathit{out}_0 = \lesf(S')\,\|\,\mathtt{22}\,\|\,\TreeSpk(r',m)$, $\mathit{out}_1$ the payout; $\mathit{extra}$ raw \\
spend type, input index & constant $\mathtt{02\,00000000}$ & script path, no annex, input 0 \\
tapleaf hash & raw & the leaf being executed \\
key version, codesep & constant $\mathtt{00\,ffffffff}$ & \\
\bottomrule
\end{tabular}
\end{center}

The message is hashed with the $\mathsf{TapSighash}$ tag and the
challenge $e$ is computed from $G_x \,\|\, G_x \,\|\, \mathsf{msg}$ under
the challenge tag. The script verifies $e_{[0,31)} \,\|\, \mathtt{00} = e$
against the witness-supplied $e_{[0,31)}$ and finally executes
\[
\langle G_x \,\|\, e_{[0,31)} \,\|\, \mathtt{01}\rangle\ \ \langle G_x\rangle\ \ \op{CHECKSIG}.
\]

\paragraph{Grinding and fee bumping.} The spender grinds the fee input's
\texttt{nSequence} until $e$ ends in $\mathtt{00}$ (so that $e + 1$ is
formed by replacing the last byte), 256 hashes on average, choosing
values that set no BIP-68 relative lock (the disable bit set) and stay
below $\mathtt{0xFFFFFFFE}$ so the transaction keeps signalling
replaceability. Any change to the transaction changes the message, so a
replacement with a higher fee must be re-ground and re-signed by the
owner; a third party bumps it only by a child spending the fee change or
an anchor. The raw $\mathit{extra}$ outputs travel as one witness element,
and Bitcoin Core's policy limits \texttt{0xC0} witness items to 80 bytes:
one P2TR change output and a P2A anchor (56 bytes) fit, two P2TR change
outputs (86 bytes) make the exit non-standard. The simulator checks
consensus limits only.
